\documentclass[11pt]{article}
\usepackage[utf8]{inputenc}
\usepackage[T1]{fontenc}
\usepackage[margin=1in]{geometry}
\usepackage{parskip}
\usepackage{amssymb,amsmath}
\usepackage{hyperref}

\begin{document}

\begin{flushright}
K\'evin Remondi\`ere\\
Independent researcher\\
Oloron-Sainte-Marie, France\\
ORCID: 0009-0008-2443-7166\\
kevin.remondiere [at] gmail.com\\[1em]
May 11, 2026
\end{flushright}

\bigskip

\noindent Editor-in-Chief\\
\emph{Journal of Number Theory}\\
Elsevier

\bigskip

\noindent Dear Editor,

I am pleased to submit for your consideration the manuscript

\smallskip
\centerline{\textbf{``Real quadratic discriminants with integer Hurwitz ratio $D^{2}/B_{2,\chi_{D}}$:}}
\centerline{\textbf{a finite list of seven''}}
\smallskip

\noindent for publication in the \emph{Journal of Number Theory}.

The paper is a short note (8 pages) addressing a question that, to the
author's best literature review, has not previously been stated and
settled in this exact form.  Specifically, for a positive fundamental
discriminant $D$ and the associated even Kronecker character
$\chi_{D}$, the rational number
\[
k(D)\;:=\;\frac{D^{2}}{B_{2,\chi_{D}}}\;=\;\frac{\pi^{2}\sqrt{D}}{L(2,\chi_{D})}\;=\;\frac{D^{2}}{24\,\zeta_{K}(-1)},\qquad K=\mathbb{Q}(\sqrt{D}),
\]
is shown to be an \emph{integer} for \emph{exactly seven} values
$D\in\{8,12,24,28,60,76,156\}$, with $k(D)\in\{32,36,48,49,75,76,117\}$
respectively, when ranging over all fundamental $D\le 10^{6}$.

The closed form for $L(2,\chi_{D})$ is the classical Hurwitz formula
(Hurwitz 1882, Iwasawa 1972, Washington 1997), and the values
$24\zeta_{K}(-1)$ are accessible via Zagier's 1976 tables.  What
appears to be new is:

\begin{itemize}
\item \emph{Explicit normalisation and classification by integrality of
$k(D)=D^{2}/B_{2,\chi_{D}}$}, rather than by integrality of
$24\zeta_{K}(-1)$ alone.  This refinement isolates a strictly smaller
set than the one ``$24\zeta_{K}(-1)\in\mathbb{Z}$'' that is classically
known to be infinite (true for all fundamental $D$).
\item \emph{The explicit finite list of seven elements and its
computational verification up to $D=10^{6}$}, together with a
heuristic density argument ($\sim D^{-3/2}$) supporting completeness.
\item \emph{The Hirzebruch--Zagier signature interpretation}, identifying
the seven discriminants as exactly those for which the arithmetic
part of the Hilbert modular surface signature divides $D^{2}$.
\item \emph{The parity-dual observation} that the imaginary-quadratic
analogue at $s=3$ produces exactly one integer-$k$ case ($d=-4$),
the simplest odd-character mirror of the seven-element real list.
\end{itemize}

The proof of the seven-element classification is short: a direct
evaluation of $B_{2,\chi_{D}}$ for each of the seven values, together
with an exhaustive sweep verified by two independent methods
(direct Bernoulli sum and \texttt{lfun}-based functional-equation
evaluation).  The verification protocol is reproduced inside the
paper for the reader's convenience.

\textbf{Manuscript length}: 8 pages, including 1 table, 1 numerical
appendix and 19 references.

\textbf{No conflict of interest.}  The author has no financial or
non-financial conflicts to disclose.

\textbf{Suggested reviewers}: scholars who have published on
$\zeta_{K}(-1)$ for real quadratic fields, or on Hilbert modular surface
signature, will be particularly well placed to evaluate this work.
Suggestions (declared without solicitation):

\begin{itemize}
\item Don Zagier (Max-Planck-Institut f\"ur Mathematik, Bonn),
\item Lawrence~C.~Washington (University of Maryland),
\item Gerard van der Geer (Universiteit van Amsterdam).
\end{itemize}

\textbf{Status}: original submission, not under review elsewhere.

I am grateful for your consideration.

\bigskip\bigskip

\noindent Sincerely,

\bigskip
\noindent K\'evin Remondi\`ere

\end{document}
